\answer{ex:03:1}{$35\to17.5\mV$ pour $g_E=g_L$; $56\to40\mV$ pour $g_E=4g_L$, une soustraction donnerait $38.5\mV$: le shunt divise $g_E$ par $3$. La fenêtre est réduite de moitié ($\tau_{\text{eff}}$ est divisée par $2$).}
\answer{ex:03:2}{$\operatorname{tr}=\frac{w_{EE}-1}{\tau_E}-\frac1{\tau_I}$, $\det=\frac{1-w_{EE}+w_{EI}w_{IE}}{\tau_E\tau_I}$; à la limite $\omega=\sqrt{\det}$, $\tau_I=20\ms$, période $73\ms$.}
\answer{ex:03:3}{$d_c=\tau_E\arccos(-1/G)/\sqrt{G^2-1}$, période $2\pi\tau_E/\sqrt{G^2-1}\in(2d_c,4d_c)$. $G=2$: $d_c=12.1\ms$, période $36\ms$; $G=5$: $d_c=3.6\ms$, période $12.8\ms$, en plein dans la bande des rythmes rapides.}
\answer{ex:03:4}{(a)~À $I_2=\beta I_1=2$ en montant et à $I_2=I_1/\beta=0.5$ en descendant: une boucle d'hystérésis de largeur $1.5$. (b)~De $\tau\ln10/(\beta-1)=2.3\tau$ par ordre de grandeur de $\varepsilon$.}
\answer{ex:03:5}{Trois intermédiaires, $d_k=5$, $10$, $15\ms$: les vetos doivent couvrir $15\ms$, à raison de $5\ms$ chacun.}
\answer{ex:03:6}{$n+M+1=3+4+1=8$ neurones. Avec deux: \nt{GABA} $=[x+y+z\ge2]$, sortie $=[x+y+z-2\,\nt{GABA}\ge1]$.}
\answer{ex:03:7}{$\binom84=70<100\le\binom94=126$: $9$ intermédiaires (théorème de Sperner). Sans connexions directes: $100$, car le rang de $\beta(J-I)$ vaut $n$.}
\answer{ex:03:8}{$h_c=\frac1{2w_{EE}}+\frac{w_{EI}w_{IE}-w_{EE}}{4w_{EE}^2}=\frac7{18}\approx0.39$; la simulation donne un changement de signe entre $h=0.38$ et $0.40$.}
